\subsection{Experiment Results}
\label{sec:experiment_res}

% \begin{table}
% \centering
% \tiny
% \begin{tabular}{ccccccc}
% \toprule
% & \textbf{Pre-train} & \textbf{B@1} & \textbf{B@2} & \textbf{METEOR} & \textbf{CIDEr} \\
% \midrule
% \multirow{2}{*}{\rotatebox{90}{Plot}} & \checkmark & 19.5 & 10.1 & 8.7 & 10.9 \\
% & \ding{55} & 20.4 & 10.7 & 9.0 & 13.1 \\
% \midrule
% \multirow{2}{*}{\rotatebox{90}{BSC}} & \checkmark & 9.1 & 3.9 & 4.3 & 7.6 \\
% & \ding{55}  & 10.1 & 4.3 & 4.4 & 8.5 \\
% \bottomrule
% \end{tabular}
% \caption{\label{tab:caption}
% The plot text generation results of \TVD.
% }
% \end{table}


% \begin{table*}
% \centering
% \tiny
% \begin{tabular}{ccccccc}
% \toprule
% \textbf{Video} & \multicolumn{2}{c}{\textbf{Text}} \\
% \midrule
% \multirow{4}{*}{\includegraphics[width=0.3\linewidth]{figure/fig_gen_1.png}} & \multirow{2}{*}{Ground Truth:} & \begin{CJK*}{UTF8}{gbsn}见吉米决定不去，她也就不再坚持。\end{CJK*} \\
%  & & When Jimmy decided not to go, she stopped insisting. \\
%  & \multirow{2}{*}{Generation Result:} & \begin{CJK*}{UTF8}{gbsn}吉米兴致不高，金没有什么反应。\end{CJK*} \\
%  & & Jimmy was not very interested and Kim did not react. \\
%  & & \\
% \multirow{4}{*}{\includegraphics[width=0.3\linewidth]{figure/fig_gen_2.png}} & \multirow{2}{*}{Ground Truth:} & \begin{CJK*}{UTF8}{gbsn}泰勒一直都没有和凯若琳联系，这让凯若琳有些生气。\end{CJK*}  \\
%  & & Taylor has not been in touch with Caroline, which makes Caroline a little angry. \\
%  & \multirow{2}{*}{Generation Result:} & \begin{CJK*}{UTF8}{gbsn}凯若琳很不高兴，她告诉艾琳娜自己还是要去找自己的母亲。\end{CJK*} \\
%  & & Caroline is very unhappy, she tells Elena that she is still going to find her mother. \\
%  & & \\
% \multirow{4}{*}{\includegraphics[width=0.3\linewidth]{figure/fig_gen_3.png}} & \multirow{2}{*}{Ground Truth:} & \begin{CJK*}{UTF8}{gbsn}老迈克看到那两人将手做出手枪的模样指向凯莉。\end{CJK*}  \\
%  & & Old Mike saw the two men point their hands like pistols at Kelly. \\
%  & \multirow{2}{*}{Generation Result:} & \begin{CJK*}{UTF8}{gbsn}那两个手下还有枪，并指向了他。\end{CJK*} \\
%  & & The two men still had guns and pointed them at him. \\
% \bottomrule
% \end{tabular}
% \caption{\label{tab:plot_generation_case}
% Case study of plot text generation in \TVD.
% }
% \end{table*}




% \setlength\tabcolsep{1.1mm}
% \begin{table*}[t]
% \centering
% \resizebox{!}{1.7cm}{
% \tiny
% \begin{tabular}{ccccccc}
% \toprule
%  & \textbf{Ref \#} & \textbf{P-T} & \textbf{BELU@1} & \textbf{BELU@2} & \textbf{METEOR} & \textbf{CIDEr} \\
% \midrule
% \multirow{2}{*}{\rotatebox{90}{Plot}} & \multirow{2}{*}1 & \checkmark & 19.5 & 10.1 & 8.7 & 10.9 \\
% & & \ding{55} & 20.4 & 10.7 & 9.0 & 13.1 \\
% \midrule
% \multirow{2}{*}{\rotatebox{90}{BSC}} & \multirow{2}{*}1 & \checkmark& \textcolor{white}{0}9.1 & \textcolor{white}{0}3.9 & 4.3 & \textcolor{white}{0}7.6 \\
% & & \ding{55}  & 10.1 & \textcolor{white}{0}4.3 & 4.4 & \textcolor{white}{0}8.5 \\
% \midrule
% \multirow{2}{*}{\rotatebox{90}{BSC}}& \multirow{2}{*}{5}  & \checkmark & 36.6 & 17.8 & 8.8 & \textcolor{white}{0}5.4 \\
% & & \ding{55}  & 32.6 & 15.6 & 8.7 & \textcolor{white}{0}5.7 \\
% \bottomrule
% \end{tabular}
% }
% \caption{Results of Plot Text Generation. ``Ref \#'' and ``P-T'' respectively represent the number of golden references and whether pre-training on BSCs is enabled.} 
% \label{tab:plot_main_tab}
% \end{table*}

% In the experiment, we mainly focus on the data of text (plot and BSC), image, and video modalities, and the data statistics as shown in Tab.~\ref{tab:merge_tab} (a).

\noindent\textbf{Genre Classification.} From Tab.~\ref{tab:ablation_genre_cls}, we confirm that multimodal data does help in the Genre Classification task, as the configuration where all modalities are involved clearly outperformed the ablation setups. 
Moreover, we noticed that text signals, especially the plot text, substantially contribute more than visual ones. Such a trend is also observed in uni-modal experiments, which are presented in the appendix due to the length limit.
% Appendix~\ref{sec:app_unimodal_genre_cls} due to the length limit.

\begin{wraptable}{r}{0.5\textwidth}
\centering
\scriptsize
\begin{tabular}{lcccc}
\toprule
\textbf{Modality} & \textbf{Mirco-F1} & \textbf{Drop} & \textbf{Marco-F1} & \textbf{Drop} \\
\midrule
All & 95.3 & - & 82.8 & - \\
\textit{w/o} Plot & 78.7 & 16.6 & 75.4 & 7.4 \\
\textit{w/o} BSC & 90.8 & \textcolor{white}{0}4.5 & 80.2 & 2.6 \\
\textit{w/o} Video & 90.9 & \textcolor{white}{0}4.4 & 80.0 & 2.8 \\
\textit{w/o} Image & 92.5 & \textcolor{white}{0}2.8 & 80.8 & 2.0 \\
\bottomrule
\end{tabular}
\caption{Results of Genre Classification. We list the relative F1 drop in each ablation setup.}
\label{tab:ablation_genre_cls}
\end{wraptable}


% It proves that manually generated text is more conducive to judging TV drama genre tags than analog signals at this stage, and plot texts effectively highlight genre characteristics.

One interesting observation is that, while including more modalities substantially improves the classification performance (both Micro and Macro), meanwhile it consistently aggregates the gap between Micro and Macro scores. For instance, without plot text, Micro F1 is higher than Macro F1 by 3.3, yet using plot text widens this difference to 12.5. Larger Micro-Macro gaps indicate a more significant imbalance between the model performance in various categories. Further investigations revealed that, frequent genres (e.g., \textit{Fantasy Romance}, \textit{High-IQ Crime}, and \textit{Situation Comedy}) which are tagged on many episodes (the model thus performs strongly even with fewer modalities) tend to benefit more than the less popular ones (e.g., \textit{Documentary and Adventure}). Therefore, we conclude that the improving effects of leveraging modalities can be biased towards frequent tags, which, to our knowledge, is the first of such findings in the field of cinema/TV multimodal learning.

% Compared with the Micro metric which pays more attention to the proportion of correctly tagged samples in all tagged samples, the Macro metric in Tab.~\ref{tab:merge_tab} (b) is significantly lower than the Micro one in all settings, which indicates that in a considerable number of tags, the existing models are cannot make correct decisions. In particular, the removal of the plot text brings the Macro and Micro metrics close while drastically reducing the performance, which shows that the modal signals except the plot text perform similarly in all categories, while the plot text does not. It will help in the classification of some tags, but some cannot (i.e., widen the performance gap between different tags).


\begin{table*}[t]
\centering
\scriptsize
\begin{tabular}{ccc|ccc|ccc|c}
\toprule
\multirow{2}{*}{\textbf{Modalities}} & \multirow{2}{*}{\textbf{Pre-Train}} & \multirow{2}{*}{\textbf{Fine-Tune}} & \multicolumn{3}{c|}{\textbf{From Text}} & \multicolumn{3}{c|}{\textbf{To Text}} & \multirow{2}{*}{\begin{tabular}[x]{@{}c@{}}\textbf{Matching}\\\textbf{Accuracy}\end{tabular}} \\
\cmidrule{4-9}
 &  &  & \textbf{Hit@1} & \textbf{Hit@5} & \textbf{Hit@10} & \textbf{Hit@1} & \textbf{Hit@5} & \textbf{Hit@10} &  \\
\midrule
\multirow{3}{*}{Image-Text} & \checkmark & \ding{55} & 0.3 & 2.2 & \textcolor{white}{0}4.0 & 0.3 & 1.9 & \textcolor{white}{0}4.5 & 50.3  \\
 & \ding{55} & \checkmark & 0.4 & 2.8 & \textcolor{white}{0}6.1 & 0.5 & 3.1 & \textcolor{white}{0}5.9 & 59.1 \\
 & \checkmark & \checkmark & 0.5 & 3.4 & \textcolor{white}{0}7.0 & 0.7 & 4.8 & \textcolor{white}{0}7.5 & 63.6 \\
% \midrule
% & \textbf{Pre-Training} & \textbf{Fine-Tuning} & \textbf{VR@1} & \textbf{VR@5} & \textbf{VR@10} & \textbf{TR@1} & \textbf{TR@5} & \textbf{TR@10} & \textbf{VTM Acc.} \\
\midrule
\multirow{3}{*}{Video-Text} & \checkmark & \ding{55} & 0.2 & 0.6 & \textcolor{white}{0}1.9 & 0.1 & \textcolor{white}{0}0.6 & \textcolor{white}{0}1.4 & 79.2 \\
 & \ding{55} & \checkmark & 1.0 & 6.0 & 10.9 & 1.3 & \textcolor{white}{0}6.4 & 12.1 & 81.5 \\
 & \checkmark & \checkmark & 2.2 & 9.6 & 17.5 & 2.4 & 10.7 & 18.5 & 89.1 \\
\bottomrule
\end{tabular}
\caption{\label{tab:retrieval}
Results (\%) of Cross-Modal Retrieval. 
}
\vspace{-0.1in}
\end{table*}

\noindent\textbf{Plot Retrieval.} As shown in Tab.~\ref{tab:retrieval}, on the one hand, in all setups tested, fine-tuning with plot text substantially improves Plot Retrieval performance. BSCs (used at the pre-training stage) also benefit the final result (which is expected in \cref{sec:framework}), although not as significantly as the professionally written plot text. 

On the other hand, despite being quite effective on the cognitively shallower Genre Classification task, our large-scale framework performs unsatisfactorily on more challenging Plot Retrieval tasks. To be exact, in all setups, even Hit@10 scores can only reach up to 18.5\%. As for the Matching Accuracy, considering a random guess will yield 50\%, the results, especially those of the Image-Text group, are to be improved. In particular, we notice that when setting text (i.e., plot-focused sentences) as the retrieval target, the model performance is significantly ($p<0.5$) better than configuring videos or images as the target. This suggests that the current framework tends to distinguish \TVD's text signals better than the visual ones.

% Another outstanding trend is that, while the model achieves relatively good results in terms of Matching Accuracy, the Hit\@$n$ scores are rather low. This reveals that although our framework is able to  One explanation is that performance of ITM/VTM Acc. in  This phenomenon directly indicates that the matching performance between different modalities is relatively high, but the density of semantic distribution within each modality is very large (i.e., poor separability), which indirectly leads to a poor hit rate during retrieval. At the same time, we observe that TR (i.e., retrieving the ranked target as text) performs consistently better than IR/VR (test result is significant with p-value: 0.0018 (alpha=0.05)~\citep{P18-1128}), which indicates that ranking text with visual signals can achieve better performance. We can further infer that the separability within the text modality is better than that of the visual modality.

On conventional scene-oriented corpora, previous studies reported that inputting video or key frame does not have a remarkable performance impact~\citep{keyframe2022revisiting}. However, this does not apply to plot-oriented \TVD, where the model employing video is consistently superior to the one using the image. At the same time, we have also reached the same conclusion as~\citep{lei2020tvr}, which suggests that video data comprising subtitles and images perform better in terms of retrieval performance compared to video data comprising only multiple images. For further details, please refer to the appendix. Besides, we found that improving the interpolation frequency, which consequently helps feed more information to the model, can further improve the retrieval scores  (see experiment results in the appendix). In summary, plot-oriented \TVD can access the model's ability to capture long-form relationships while previous scene-oriented datasets cannot.

% ~\cref{sec:app_number_frame}


% The results of cross-modal retrieval are listed in , which is very intuitive. Unlike the conclusion of existing multi-modal research that claims only needs to focus on specific key frames, the significant retrieval performance gap between image and video indicates that more information needs to be considered to model the complete plot in TV drama (i.e., long video clips).


\begin{wraptable}{r}{0.6\textwidth}
\centering
\scriptsize
\begin{tabular}{ccccccc}
\toprule
 & \textbf{Ref \#} & \textbf{P-T} & \textbf{BELU@1} & \textbf{BELU@2} & \textbf{METEOR} & \textbf{CIDEr} \\
\midrule
\multirow{2}{*}{\rotatebox{90}{Plot}} & \multirow{2}{*}1 & \checkmark & 18.2 & \textcolor{white}{0}9.6 & 8.0 & \textcolor{white}{0}9.0 \\
& & \ding{55} & 15.2 & \textcolor{white}{0}7.6 & 7.6 & \textcolor{white}{0}6.7 \\
\midrule
\multirow{2}{*}{\rotatebox{90}{BSC}} & \multirow{2}{*}1 & \checkmark& \textcolor{white}{0}8.6 & \textcolor{white}{0}3.9 & 3.7 & \textcolor{white}{0}8.0 \\
& & \ding{55}  & \textcolor{white}{0}9.9 & \textcolor{white}{0}4.1 & 4.7 & \textcolor{white}{0}9.1 \\
\midrule
\multirow{2}{*}{\rotatebox{90}{BSC}}& \multirow{2}{*}{5}  & \checkmark & 25.2 & 12.7 & 8.1 & \textcolor{white}{0}6.4 \\
& & \ding{55}  & 38.9 & 18.0 & 9.2 & \textcolor{white}{0}5.6 \\
\bottomrule
\end{tabular}
\caption{Results of Plot Text Generation. ``Ref \#'' and ``P-T'' respectively represent the number of golden references and whether pre-training on BSCs is enabled.} 
\label{tab:plot_main_tab}
\end{wraptable}

\noindent\textbf{Plot Text Generation.} The results in Tab.~\ref{tab:plot_main_tab} are partial counter-intuition: on the one hand, pre-training with BSCs benefits the model on Plot Retrieval, and
constantly has a positive impact on Plot Text Generation. On the other hand, surprisingly, pre-training on BSCs brings a negative impact on all automatic evaluation scores in BSC generation. 
Even if we expanded the golden reference set to the top 5 most upvoted BSCs, the gap (BELU\@1, 2) brought by pre-training become even greater.
We have reason to suspect that the inconsistent distribution between BSC and plot text is a potential cause. We train the framework to generate BSC, rather than using the BSC with the most votes in each clip as a golden reference. We left investigation into this matter as an important future work.

Moreover, eight different text entropy calculators~\citep{entropy2022acl} suggested that the distribution of plot text is more complex than BSCs (see statistics in the appendix). Considering BSCs are also more than plot text in terms of the total number, writing BSCs is intuitively less challenging than producing text plots. Nonetheless, empirically the generation scores of plot text are higher than those of BSCs (with one reference). In the future, we will dive deeper into \TVD for explanations. 

% in~\cref{sec:bsc_plot}


\setlength\tabcolsep{1.5mm}
\begin{table*}[t]
\centering
\tiny
\begin{tabular}{cll}
\toprule
\textbf{Video} & \multicolumn{2}{c}{\textbf{Text}} \\
\midrule
\multirow{4}{*}{\includegraphics[width=0.35\linewidth]{figure/fig_gen_1.png}} & \multirow{2}{*}{Ground truth:} & \begin{CJK*}{UTF8}{gbsn}见吉米决定不去，她也就不再坚持。\end{CJK*} \\
 & & When Jimmy decided not to go, she stopped insisting. \\
 & \multirow{2}{*}{Generation:} & \begin{CJK*}{UTF8}{gbsn}吉米兴致不高，金没有什么反应。\end{CJK*} \\
 & & Jimmy was not very interested and Kim did not react. \\
 & & \\
\multirow{4}{*}{\includegraphics[width=0.35\linewidth]{figure/fig_gen_2.png}} & \multirow{2}{*}{Ground truth:} & \begin{CJK*}{UTF8}{gbsn}泰勒一直都没有和凯若琳联系，这让凯若琳有些生气。\end{CJK*}  \\
 & & Taylor has not been in touch with Caroline, which makes Caroline a little angry. \\
 & \multirow{2}{*}{Generation:} & \begin{CJK*}{UTF8}{gbsn}凯若琳很不高兴，她告诉艾琳娜自己还是要去找自己的母亲。\end{CJK*} \\
 & & Caroline is very unhappy, she tells Elena that she is still going to find her mother. \\
 & & \\
\multirow{4}{*}{\includegraphics[width=0.35\linewidth]{figure/fig_gen_3.png}} & \multirow{2}{*}{Ground truth:} & \begin{CJK*}{UTF8}{gbsn}老迈克看到那两人将手做出手枪的模样指向凯莉。\end{CJK*}  \\
 & & Old Mike saw the two men point their hands like pistols at Kelly. \\
 & \multirow{2}{*}{Generation:} & \begin{CJK*}{UTF8}{gbsn}那两个手下还有枪，并指向了他。\end{CJK*} \\
 & & The two men still had guns and pointed them at him. \\
\bottomrule
\end{tabular}
\caption{\label{tab:plot_generation_case}
Examples of generated plot text.
}
\vspace{-0.2in}
\end{table*}
 
Lastly, we listed some Plot Text Generation examples in Tab.~\ref{tab:plot_generation_case} (with more in the appendix).
% Tab.~\ref{tab:add_plot_generation_case}
These cases reveal that, trained on \TVD, our framework is able to integrate situational context and mental states to some degree. 
However, it also makes mistakes, e.g., mistaking one role for another,  misunderstanding the relationships between roles, or misinterpreting the intentions of certain actions. 
In summary, using \TVD we show there is still much room for improvement in multi-modal learning on cinema/TV data.